\par\noindent\textbf{(1)} 沿轨迹 z=vt，相距时间Δt的两处辐射在观察方向θ上的相位差为 ΔΦ=ωΔt−k(ω)vΔt cosθ。远场相干增强要求 ΔΦ=0，即 cosθ\_C=ω/[k(ω)v]=c/[n(ω)v]；这与以相速度 c/n(ω) 构造的波前包络锥一致。
\par\noindent\textbf{(2)} 锥角有实解必须 β n(ω)>1，β=v/c。若n随频率变化，只能在满足此不等式的频带存在Cherenkov辐射；不同ω对应不同锥角，形成色散角谱。对相同能量的粒子，阈速随材料折射率升高而降低。
\par\noindent\textbf{(3)} Frank–Tamm经典谱在SI单位中可写 d²W/(dx dω)=q²ω/(4πε₀c²)[1−c²/(v²n(ω)²)]，仅在βn(ω)>1时取非零，其他区域为零。前因子量纲为能量/长度/角频率，方括号来自横向辐射场投影及锥角的sin²θ\_C。
\par\noindent\textbf{(4)} 如果错误地把无色散常数n延伸到ω→∞，积分∫ωdω将出现紫外发散。真实介质在高频通常 n(ω)→1，粒子 β<1 时失去超光相速条件，材料吸收和量子反冲也限制经典理论。对有限轨道长度，delta型相位匹配会展宽而不能直接沿用无限轨迹极限。
